% FILE: 06_contribution_to_thesis.tex
% PROJECT: phd-thesis
% MANUFACTURED: 2026-06-01
% AUTHORITY: ledgers/PROJECT_INDEX.yaml, checkpoints/PI_RESEARCH_PROGRAM_ALIVE_001.md,
%            receipts/rec_ebitda_rework_001.json, phd-thesis/README.md
% DO NOT CLAIM: See ledgers/DO_NOT_CLAIM_LEDGER.md

\section{Contribution to the Dissertation}
\label{sec:phd-thesis-contribution}

\subsection{Contribution to Prediction vs.~Coordination}
\label{sec:phd-thesis-contribution-prediction}

The \texttt{phd-thesis} workspace contributes to the prediction-versus-coordination argument
by being, itself, an instance of the coordination regime. A prediction-only system would
produce dissertation content from a language model applied to the research corpus without
receipts, without a named manufacturing pipeline, and without an ALIVE verification criterion.
The \texttt{phd-thesis} workspace rejects this approach by design.

The workspace enforces the coordination regime through three mechanisms. First, the
DO\_NOT\_CLAIM\_LEDGER prohibits claims without source artifact backing --- this is the
workspace-level instantiation of the named-witness requirement that distinguishes conformance
claims from prediction-derived assertions. Second, the WORKFLOW\_RECEIPT.yaml binds every
manufacturing event to a specific toolchain, branch, and pre-manufacture commit --- this is
the workspace-level instantiation of the receipt doctrine. Third, the ALIVE criteria require
PDF compilation and hash recording before the workspace can be released --- this is the
workspace-level instantiation of the proof gate.

The dissertation therefore contributes to the prediction-versus-coordination argument not only
as a document describing the argument but as an artifact that instantiates the coordination
architecture. The workspace's PARTIAL status at time of writing is an honest record of the
manufacturing state; it does not claim to be more complete than it is. This honesty is itself
a property of the coordination regime: the event log (the manufacturing provenance) tells the
truth about what has happened, regardless of what the system might prefer to declare.

The research lineage --- 2016 LM experiment $\to$ local text imitation does not conserve
consequence $\to$ enterprise process gap $\to$ Chatman Equation $A = \mu(O^*)$ --- is the
conceptual arc that this workspace embodies. The workspace is proof that the arc is not merely
theoretical: a real manufacturing system, with real receipts and real proof gates, can be built
to enforce the coordination discipline at the artifact level.

\subsection{Contribution to Process Evidence}
\label{sec:phd-thesis-contribution-evidence}

The workspace contributes to the process evidence strand of the dissertation by being the
authoritative index of the foundry's process evidence surfaces. The
\texttt{ledgers/PROJECT\_INDEX.yaml} maps all 34 projects to their research surfaces, which
include: \texttt{process-evidence}, \texttt{conformance-checking}, \texttt{mining-authority},
\texttt{replay-authority}, \texttt{lifecycle-governance}, \texttt{receipt-chain},
\texttt{OCEL-v2}, \texttt{type-law}, \texttt{witness-lattice}, \texttt{proof-gate},
\texttt{audit-suite}, and \texttt{gap-analysis}.

This mapping is the dissertation's contribution to process evidence as an architectural
concept: it demonstrates that process evidence is not a single artifact but a layered system
of interlocking surfaces. The type law (Evidence$\langle T, \text{State}, W\rangle$) provides
the carrier; the receipt chain provides the provenance; the conformance checking surface
(wasm4pm) provides the computational authority; the adversarial surface provides the
negative-test coverage required by the Van der Aalst Constitution; and the gap registry
provides the honest record of what is not yet known.

The five canonical receipts in \texttt{receipts/} (EBITDA rework through residual standard)
are the dissertation's primary process evidence artifacts in the M\&A domain. Each receipt
carries a process model hash, an event log hash, wasm4pm-computed fitness and precision
scores, and an ed25519 validator signature. The EBITDA receipt
(\texttt{rec\_ebitda\_rework\_001.json}) alone documents fitness 0.982, precision 0.945,
throughput 4.12 days, and a projected EBITDA impact of \$1,250,000 --- all derived from
replay against a declared process model, not from a prediction.

\subsection{Contribution to Manufacturing Architecture}
\label{sec:phd-thesis-contribution-manufacturing}

The workspace contributes to the manufacturing architecture strand by demonstrating that a
dissertation --- traditionally the least structured artifact in an academic research program
--- can be manufactured under the same receipted, proof-gated discipline as a software
artifact. This is an AUTHOR THESIS claim, grounded in the structure of the workspace itself
rather than in a prior publication.

The manufacturing architecture instantiated by this workspace consists of: a branch-scoped
manufacturing run (\texttt{phd-thesis-corpus-manufacture-001}), a verified toolchain
recorded in WORKFLOW\_RECEIPT.yaml, a project index manufactured by the corpus cartographer
phase, a claim prohibition ledger enforced across all artifacts, and ALIVE criteria that gate
release on PDF compilation and hash recording. This is the same architecture --- adapted to
TeX compilation rather than Rust compilation --- as the manufacturing pipelines used by
wasm4pm, Blue River Dam, and the ggen ecosystem.

The PROJECT\_INDEX.yaml, which enumerates 34 projects with thesis roles, is itself a
manufactured artifact: it was produced 2026-06-01 as the output of a corpus cartographer
phase that scanned the \texttt{process-intelligence} foundry and classified each project's
evidentiary contribution to the dissertation. This manufacturing phase is receipted in the
WORKFLOW\_RECEIPT.yaml (\texttt{phase: SETUP\_COMPLETE}) and will be followed by chapter
manufacture and PDF compilation phases, each of which will produce its own provenance record.

\subsection{Contribution to Capital-Flow Infrastructure}
\label{sec:phd-thesis-contribution-capital}

The workspace contributes to the capital-flow infrastructure strand through the M\&A claim
surfaces documented in the PROJECT\_INDEX.yaml. The \texttt{ma} project (42 M\&A claim files,
board claim taxonomy, slide-to-receipt map, reverse Porter Five, hostile takeover receipts)
and the \texttt{receipts} project (5 canonical JSON receipts with validator signatures) are
the primary capital-flow artifacts in the dissertation corpus.

The thesis lineage records the capital-flow arc as: receipts and replay $\to$ ggen/Open
Ontologies $\to$ clap-noun-verb command grammar $\to$ wasm4pm/process-evidence $\to$
Post-Cyberpunk PCP $\to$ AI XYNZ $\to$ capital-flow/settlement/DAO $\to$ industry-complete
architecture. The \texttt{phd-thesis} workspace is the point at which all of these strands
converge into a single, receipted, proof-gated corpus.

The capital-flow contribution is not a claim about deployment or commercialization --- the
DO\_NOT\_CLAIM\_LEDGER explicitly prohibits production deployment claims without receipt. It
is a claim about the admissibility architecture: that process evidence, when manufactured
under the receipt and replay doctrine, produces claims that are board-admissible and
independently verifiable. The five canonical receipts in the \texttt{receipts/} directory
are the physical instantiation of this admissibility --- each one a manufacturable, hashable,
signable unit of capital-flow process evidence.
